\chapter{A Brief History of Risk Assessment}
\label{chap:A Brief History of Risk Assessment}

%---------------------------- SECTION ----------------------------
\section{Why History Matters to a Compliance Professional}
\paragraph{It might be tempting to skip a chapter on history. After all, the pressures facing compliance and legal professionals today are immediate and practical — there are frameworks to implement, regulators to satisfy, and risks to manage right now. What does the past have to offer?}
\paragraph{Quite a lot, as it turns out.}
\paragraph{Understanding how risk assessment evolved — what drove its development, what failures exposed its absence, and what regulatory responses followed — gives context to the frameworks, obligations, and expectations you work with every day. Many of the requirements that compliance professionals now take for granted did not emerge from abstract theory. They emerged from real events, real failures, and real harm. Knowing that history helps you understand not just \textit{what} the rules are, but \textit{why} they exist — and that understanding makes you considerably more effective at applying and explaining them,}
\paragraph{It also serves as a useful reminder that the field continues to evolve. The risk landscape facing organisations today looks very different from the one that existed twenty years ago, and twenty years from now it will look different again. A profession that understands its own history is better placed to anticipate and adapt to what comes next.}

%---------------------------- SECTION ----------------------------
\section{The Earliest Roots — Risk as Old as Human Decision-Making}
\paragraph{Whilst the formal discipline of risk assessment is relatively modern, the underlying impulse — the desire to anticipate what might go wrong and prepare for it — is as old as human civilisation itself.}
\paragraph{Ancient merchants calculating the odds of a ship returning safely from a trading voyage, agricultural societies developing grain storage practices to guard against poor harvests, and early builders applying rules of thumb to ensure structures would not collapse — all of these represent, in their own way, a form of risk thinking.}
\paragraph{The formalisation of probability theory in the seventeenth century, through the work of mathematicians such as Blaise Pascal and Pierre de Fermat, gave risk thinking a mathematical foundation. The emergence of insurance markets — particularly in London's Lloyd's Coffee House in the late seventeenth century — demonstrated that risk could be quantified, priced, and transferred. These were significant intellectual and commercial developments, but they remained largely confined to financial and actuarial contexts for some time.}
\paragraph{The journey from probabilistic mathematics to the structured, governance-oriented risk assessment practices of the modern compliance world was a long one — and it was driven, more than anything else, by a series of catastrophic events that exposed the consequences of inadequate risk management.}

%---------------------------- SECTION ----------------------------
\section{The Industrial Era — When Risk Became a Public Concern}
\paragraph{The industrial revolution transformed the scale and nature of risk. Factories, railways, mines, and chemical plants introduced hazards that affected not just those who chose to take on commercial risk, but workers, communities, and the wider public — people who had no meaningful say in the matter.}
\paragraph{The response, over time, was \textbf{regulation}. Governments began to impose requirements on organisations to identify and manage the risks their operations created. Early factory legislation in the United Kingdom — beginning with the Health and Morals of Apprentices Act of 1802 and developing through successive Factory Acts throughout the nineteenth century — represents one of the earliest examples of risk-related regulatory obligation. These were crude by modern standards, but they established a principle that would prove enduring: \textbf{organisations have a duty to identify and manage risks to those affected by their activities, and the state has a role in enforcing that duty}.}
\paragraph{Similar patterns emerged across industrialising nations throughout the nineteenth and early twentieth centuries, with regulation gradually expanding to cover workplace safety, product liability, environmental protection, and financial conduct.}

%---------------------------- SECTION ----------------------------
\section{The Mid-20th Century — The Birth of Formal Risk Assessment}

\paragraph{The formal discipline of risk assessment as we would recognise it today began to take shape in the mid-twentieth century, driven largely by two distinct developments: the growth of the nuclear and chemical industries, and the increasing sophistication of financial markets.}

\subsection{The Nuclear and Chemical Industries}

\paragraph{The development of nuclear technology in the 1940s and 1950s introduced a category of risk that was genuinely unprecedented in its potential scale. The consequences of a serious nuclear accident were so severe, and so difficult to reverse, that a fundamentally new approach to risk management was required — one that went far beyond the reactive, incident-driven approaches of the industrial era.}
\paragraph{Out of this necessity came some of the foundational techniques of modern risk assessment. \textbf{Fault Tree Analysis}, developed at Bell Laboratories in the early 1960s and subsequently adopted by the nuclear industry, provided a systematic method for tracing the pathways through which a failure could occur. \textbf{Failure Mode and Effects Analysis (FMEA)}, originally developed by the US military in the 1940s, offered a structured approach to identifying how individual components might fail and what the downstream consequences would be}
\paragraph{The chemical industry made equally significant contributions. A series of major industrial accidents throughout the 1970s and 1980s — most devastatingly the Union Carbide disaster in Bhopal, India in 1984, in which a gas leak at a pesticide plant killed thousands of people and injured hundreds of thousands more — demonstrated the catastrophic consequences of inadequate process safety management. The regulatory and industry response to Bhopal and similar events accelerated the development of structured hazard analysis methodologies, most notably \textbf{HAZOP (Hazard and Operability Study)}, and drove the introduction of mandatory risk assessment requirements in chemical and process industries across many jurisdictions}

\subsection{The Financial Sector}

\paragraph{The financial sector was developing its own risk thinking in parallel. The growth of derivatives markets, the increasing complexity of financial instruments, and the globalisation of capital flows through the latter half of the twentieth century created risks that required quantitative modelling and systematic management.}
\paragraph{The 1970s and 1980s saw the emergence of increasingly sophisticated financial risk models. The development of \textbf{Value at Risk (VaR)} as a standard measure of market risk, and its adoption by major financial institutions and regulators from the 1990s onwards, marked a significant step in the formalisation of financial risk assessment. The Basle Accord of 1988 — the forerunner of the Basel framework that continues to shape bank capital regulation today — represented one of the first major international regulatory frameworks built explicitly around the concept of quantified risk.}

%---------------------------- SECTION ----------------------------
\section{The Late 20th Century — Regulatory Frameworks Take Shape}
\paragraph{By the 1980s and 1990s, risk assessment was becoming embedded in regulatory frameworks across a widening range of sectors. Several developments from this period are worth highlighting, as their legacy remains directly relevant to compliance professionals today.}

\subsection{Health and Safety Legislation}

\paragraph{In the United Kingdom, the \textbf{Health and Safety at Work Act 1974} had established a broad duty on employers to ensure, so far as reasonably practicable, the health, safety, and welfare of employees. The \textbf{Management of Health and Safety at Work Regulations 1992} — introduced to implement a European Community directive — made the requirement to conduct formal, documented risk assessments explicit and enforceable. This was a watershed moment: risk assessment was no longer merely good practice. It was a legal obligation with consequences for non-compliance.}

\paragraph{Similar legislative developments were taking place across Europe, North America, and beyond, establishing risk assessment as a standard element of regulatory compliance in occupational health and safety.}

\subsection{Environmental Regulation}

\paragraph{Growing awareness of environmental risk — driven by events such as the Exxon Valdez oil spill in 1989 and a series of major industrial pollution incidents — led to the development of environmental risk assessment frameworks and requirements. The principle that organisations should systematically identify and manage environmental risks, and that regulators could hold them accountable for failing to do so, became increasingly embedded in environmental law across many jurisdictions.}

\subsection{Financial Services Regulation}

\paragraph{The financial services sector saw an intensification of risk-based regulatory oversight during the 1990s. Regulators began to move away from prescriptive rule-based approaches — telling firms exactly what they must and must not do — towards \textbf{"risk-based supervision"}; an approach in which regulatory scrutiny was calibrated to the level of risk a firm posed, and in which firms were expected to demonstrate robust internal risk management capability.}
\paragraph{This shift had profound implications for compliance professionals. It meant that demonstrating how you managed risk became as important as demonstrating what you had done. The quality and credibility of your risk assessment process became, in itself, a subject of regulatory interest.}

%---------------------------- SECTION ----------------------------
\section{The 21st Century — Crisis, Complexity and a Broadening Agenda}

\paragraph{If the latter decades of the twentieth century saw risk assessment become embedded in regulatory frameworks, the early decades of the twenty-first century tested those frameworks severely — and, in many cases, found them wanting.}

\subsection{The 2008 Financial Crisis}
\paragraph{No event has done more to shape contemporary thinking about risk management in financial services — and, arguably, across the broader compliance landscape — than the global financial crisis of 2007 to 2009.}
\paragraph{The crisis exposed a series of deeply troubling failures in risk assessment and risk management: financial institutions that had built complex, opaque risk exposures that their own senior management did not fully understand; risk models that had systematically underestimated the possibility of correlated failures across markets; governance structures that had failed to challenge or constrain risk-taking; and regulatory frameworks that had placed excessive confidence in the internal risk models of the very institutions they were meant to oversee.}
\paragraph{The post-crisis regulatory response was sweeping. The \textbf{Basel III framework} significantly tightened capital and liquidity requirements for banks and introduced new expectations around risk governance. In the United Kingdom, the \textbf{Senior Managers and Certification Regime (SM\&CR)} — developed in response to the parliamentary commission on banking standards — introduced personal accountability for senior individuals at regulated firms, making it harder for executives to distance themselves from failures of risk oversight. In the United States, the \textbf{Dodd-Frank Act} introduced a range of new risk-related requirements and significantly expanded the regulatory perimeter.}
\paragraph{Beyond the specific regulatory changes, the financial crisis fundamentally shifted the culture of expectation around risk management. The idea that sophisticated quantitative models were sufficient, that risk management was primarily a technical function, and that senior management could rely on risk professionals to keep things in check without deep personal engagement — all of these assumptions were challenged and, in many cases, discredited.}

\subsection{Data Protection and Cybersecurity}
\paragraph{The rapid expansion of digital technology and data-driven business models through the 2000s and 2010s created an entirely new category of risk that existing frameworks were not designed to address. The introduction of the \textbf{General Data Protection Regulation (GDPR)} in the European Union in 2018 — and the wave of data protection legislation it influenced globally — brought risk assessment explicitly into the field of data governance. The requirement to conduct \textbf{Data Protection Impact Assessments (DPIAs)} for high-risk processing activities is, at its core, a risk assessment obligation translated into data protection law.}
\paragraph{Simultaneously, the growing threat of cyber attack created new urgency around cybersecurity risk assessment. High-profile breaches affecting major organisations across every sector demonstrated that cyber risk was neither a niche technology concern nor a theoretical one. Regulators in financial services, healthcare, critical infrastructure, and elsewhere responded with increasingly specific requirements around cyber risk identification, assessment, and management.}

\subsection{ESG and Non-Financial Risk}
\paragraph{Perhaps the most significant recent evolution in the scope of risk assessment has been the growing attention paid to \textbf{Environmental, Social, and Governance (ESG)} risks. Driven by a combination of investor pressure, regulatory development, and genuine shifts in societal expectations, organisations across every sector are now expected to identify, assess, and report on risks related to climate change, human rights in supply chains, diversity and inclusion, and governance quality}
\paragraph{This represents a meaningful expansion of the risk assessment agenda — one that is still developing rapidly and that is likely to generate significant new compliance obligations in the years ahead. We will explore this in more detail in Chapter~\ref{chap:Environmental and Sustainability Risk Assessment}}

%---------------------------- SECTION ----------------------------
\section{From Compliance Teams to Risk Assessors — How We Got Here}
\paragraph{Looking back across this history, a clear pattern emerges. Time and again, significant events — industrial accidents, financial crises, environmental disasters, data breaches — have exposed the inadequacy of existing risk management practices, prompted regulatory responses, and ultimately created new obligations and expectations for the organisations and professionals responsible for compliance.}
\paragraph{This pattern explains why compliance and legal teams find themselves so deeply involved in risk assessment today. It is not that risk assessment was invented for compliance purposes. It is that the consequences of inadequate risk management have, repeatedly and consistently, been addressed through regulation — and regulation is the domain of compliance.}
\paragraph{The compliance professional's relationship with risk assessment is therefore not incidental. It is the product of decades of hard lessons, regulatory evolution, and the gradual recognition that organisations need people who can think systematically about what could go wrong, ensure that appropriate processes are in place to manage it, and demonstrate credibly to regulators and stakeholders that this is being done properly.}
\paragraph{That is the role you occupy. And understanding the history of how it came to be is the beginning of performing it well.}

%---------------------------- SECTION ----------------------------
\section{Chapter Summary}
\paragraph{Before moving on, let us anchor the key ideas from this chapter:}

\paragraph{Key takeaways:}

\sloppy
\begin{itemize}
  \bitem{Pre-industrial}{Informal risk thinking; emergence of probability theory and insurance markets.}
  \bitem{Industrial revolution}{Risk as a public concern; early regulatory responses to workplace and public harm.}
  \bitem{Mid-20th century}{Formal risk assessment techniques developed in nuclear, chemical, and financial sectors.}
  \bitem{Late 20th century}{Risk assessment embedded in regulatory frameworks — health \& safety, environment, financial services.}
  \bitem{2008 financial crisis}{Systemic risk management failures; sweeping regulatory reform; personal accountability introduced.}
  \bitem{2010s–present}{Cyber, data protection, and ESG risks expand the risk assessment agenda significantly.}
  \bitem{Takeaways}{In this chapter we learned:
  \begin{itemize}
          \item{Formal risk assessment grew not from theory, but from the consequences of real failures and the regulatory responses they triggered.}
          \item{Many of the compliance obligations you work with today have their roots in specific, often catastrophic, events.}
          \item{The shift to risk-based supervision means that how you manage risk is as important to regulators as what you do.}
          \item{The agenda continues to expand — cyber, data protection, and ESG are reshaping the risk landscape in real time.}
          \item{Compliance professionals occupy their central role in risk assessment for good historical reasons — understanding those reasons strengthens the case for doing the work well.}
        \end{itemize}
    }
\end{itemize}
\fussy

\barquote{In Chapter~\ref{chap:Types of Risk}, we will take stock of the modern risk landscape by exploring the full taxonomy of risk types that organisations face today — laying the groundwork for the more detailed assessment and management techniques that follow in Parts II and III.}